\chapter{实验配置与结果索引}
\label{app:experiments}

\section{已发表工作与新增实验的来源}

第三章结果图来自作者保存的已发表论文工程，原始图文件保持不变。图3-1依据模型定义重新绘制，其余图组引用已发表工作；空间动作快照只进行组合排版。原图能够支持相应曲线与空间现象的讨论，但由于缺少逐次原始运行数据，本文不据此重新计算运行间标准差或置信区间。

工作2采用独立建立的仿真程序，并以第四章给出的模型定义为准。新增实验与工作1原图保持来源分离，不将重建环境中的结果称为对已发表曲线的数值复现。完整代码、配置、逐次输出与复现说明保存在项目仓库及随附研究材料中，本附录只保留理解论文实验所需的配置与冻结参数。

\section{工作2实验配置}

表\ref{tab:app-development}给出工作2的主要实验设置。训练用于参数搜索，独立验证用于候选选择，最终留出测试在控制器冻结后执行。不同阶段使用分离的环境种子，最终测试对每个方法—条件使用30次独立运行；主要统计比较的重复单位均为完整独立运行。

\begin{table}[!htbp]
\centering
\caption{工作2主要实验配置}
\label{tab:app-development}
\begin{tabular}{ll}
\toprule
配置项 & 数值或设置 \\
\midrule
网格与默认感知范围 & $30\times30$周期网格，$M=2$ \\
状态与动作 & 二值邻域声誉状态，合作与背叛 \\
个体学习参数 & $\eta=0.8$，$\gamma=0.9$ \\
探索调度 & 初值0.5，每步乘0.99，下限0.01 \\
综合奖励中的收益权重 & $w_P=1$ \\
控制器最大附加强度 & $\kappa_{\max}=0.5$ \\
训练环境 & $r\in\{4.4,4.8\}$，$\rho\in\{0,0.1\}$ \\
基础异常消息 & 持续报告奖励1与背叛动作 \\
单个训练候选时长 & 10000步 \\
每次独立搜索 & 24候选，12代，每代5个精英 \\
独立优化初始化数 & 每类结构3次 \\
验证与最终测试时长 & 100000步 \\
最终测试重复 & 每个方法--条件30次独立运行 \\
末段评价窗口 & 运行时长的最后20\% \\
\bottomrule
\end{tabular}
\par\smallskip{\zihao{-5}\rmfamily 来源：作者，本文模型定义与实验配置，2026年。\par}
\end{table}

首轮四维表示、固定规则强度比较以及未选中训练输出见附录\ref{app:development-history}。这些开发结果用于说明模型形成与选型过程，不并入冻结后的最终测试样本；更细的运行配置和随机种子记录保存在实验仓库中。

\section{冻结的控制器参数}

最终比较中的四类学习结构均在完成独立优化后，通过共同验证环境选择一组冻结参数。完整模型、固定合作评分门控、仅学习评分以及最高奖励参考下学习门控分别对应第四、五章讨论的四种结构。表\ref{tab:app-frozen-parameters}按第四章记号列出冻结参数；实际运行使用归档配置中的完整精度数值，表格中的舍入值仅用于论文展示。

\begin{table}[!htbp]
\centering\small
\caption{共同验证选定的控制器参数；展示值保留六位小数}
\label{tab:app-frozen-parameters}
\begin{tabular}{lrrrr}
\toprule
参数 & 完整模型 & 固定评分/门控 & 学习评分/固定门控 & 最高奖励/门控 \\
\midrule
$w_1$ & 0.879566 & 1.000000 & 0.459648 & 1.000000 \\
$w_2$ & -0.120373 & 0.000000 & 0.101663 & 0.000000 \\
$w_3$ & 1.366790 & 0.000000 & 1.957342 & 0.000000 \\
$w_4$ & -0.990653 & 0.000000 & -1.133374 & 0.000000 \\
$w_5$ & 2.132541 & 2.000000 & 2.296933 & 2.000000 \\
$v_1$ & 1.440104 & 1.624851 & 0.000000 & -0.914830 \\
$v_2$ & -2.549781 & -2.940219 & 0.000000 & -0.124759 \\
$v_3$ & 0.077607 & 6.258113 & 0.000000 & -2.206454 \\
$v_4$ & 1.817067 & -0.478881 & 0.000000 & -1.008254 \\
$v_5$ & 0.473214 & -1.265724 & 0.000000 & 4.510742 \\
$v_u$ & 1.192334 & -0.429508 & 0.000000 & -2.458004 \\
$b$ & -0.578214 & 2.579153 & -0.405465 & -2.290863 \\
\bottomrule
\end{tabular}
\par\smallskip{\zihao{-5}\rmfamily 来源：作者，本文模型开发与仿真实验，2026年。\par}
\end{table}


四类控制器的最大附加强度均为0.5。固定合作评分结构的参考排序保持合作优先形式；仅学习评分结构的门控恒为0.4；最高奖励结构直接按报告奖励选择参考对象。固定合作规则在主要比较中取$\kappa=0.2$，幅度校准对照取$\kappa=0.027704459701337496$，二者在最终测试中均保持冻结。

\section{补充结果索引}

开发过程中的表示修订、简单规则强度比较以及全部训练输出集中列于附录\ref{app:development-history}。冻结后的最终测试补充结果，包括完整末段表、全部运行分布、主要配对区间和补充空间快照，列于附录\ref{app:final-supplement}。方法开发、最终测试和复现工程材料分别保存，正文只保留支撑研究问题与结论所需的实验信息。